\honest{But There's Still A Chance, Right?}{Eliezer Yudkowsky}{2008}

Years ago a man told me he did not believe in evolution. I told him we can read the genes: humans and chimpanzees share 98\% of their DNA. He said the similarity might be coincidence. I said the odds against that were ``something like two to the power of seven hundred and fifty million to one.'' He said, ``But there's still a chance, right?''

I admit that I do not remember where I got that figure, and that anyone who states odds like that is wrong far more often than the odds allow. The estimate of shared DNA was also revised to 95\%, and that figure ``may even apply only to the 30,000 known genes,'' in which case my odds are off by a ``meta-order of magnitude.'' \nb{The 95\% estimate (Britten, 2002) came from genomic DNA, repeated and non-repeated sequence alike, not from genes, and the 2005 comparison of the two whole genomes found the same kinds of difference across the genome. So that case does not arise.} Still, I find his reply funny.

It taught me, I say, ``several insights into the laws of thought as seen by the unenlightened ones.'' The first is that intuition treats ``No chance'' and ``A very tiny chance'' as different in kind. \nb{Kahneman and Tversky had described this in 1979, in prospect theory.} We can write down numbers that small but cannot feel them; a feeling that small ``doesn't fire enough neurons.'' I give no source for that. ``This is why people buy lottery tickets.''

The second is that an argument that is not certain may, on this view, be ignored: at a likelihood of zero you must give up the belief, at one over a googol you may keep it. \nb{This rule is Yudkowsky's reading of the one sentence the man is reported to have said.} But if you may ignore evidence at one over a googol, I ask, why not at zero?

I learn from other people's blatant bad examples and apply the lesson to subtler cases. Here the lesson is that if you cannot ignore a likelihood of one over a googol because you want to, you cannot ignore one of 0.9 either. When you catch yourself thinking ``But you can't prove me wrong,'' remember him: if you ignore a probabilistic argument, why not ignore a proof too?
